\section{总结与延伸}

本讲的核心结论可以压缩成一句话：Codex 的 planning 能力不是单一按钮，而是一组分层机制。Plan Mode 负责在动手前把规格聊清楚；\texttt{update\_plan} 负责在执行中维护 live progress；文件化 todo 则负责把长期状态从 context window 中解耦出来。

讲者最后指出，随着模型能力增强，像 \texttt{update\_plan} 这样极简的工具已经足以让模型完成复杂的创建、状态推进和 replan。真正重要的不是工具数量，而是工具 schema 是否和模型的语义能力对齐。

\begin{importantbox}{最终 takeaway}
优秀的 agent harness 不一定要堆很多专用工具。一个足够通用、语义清晰、状态可视的工具，配合强模型，往往能覆盖更大的行为空间。Codex 的 \texttt{update\_plan} 就是这种设计的代表。
\end{importantbox}

对于我们维护课程笔记仓库，也可以借鉴这套思路：短任务用上下文计划推进，长任务用文件化 todo 固化状态；每次遇到下载、字幕、转录、编译等阻塞，都不要只停在口头记忆里，而要写回队列，形成可恢复的工作流。

\end{document}
